\documentclass[11pt]{article}

\usepackage[margin=1in]{geometry}
\usepackage{booktabs}
\usepackage{amsmath}
\usepackage{hyperref}

\title{Dual-Domain MMD Adaptation: Current Paper Status and Remaining Work}
\author{}
\date{\today}

\begin{document}
\maketitle

\section{The story so far}

\textbf{Question.} Can cheap, domain-shifted synthetic aerial imagery improve real-domain
(VisDrone) object detection via feature-level alignment (MMD) and under what conditions?

\textbf{Method.} A YOLOv10 detector extended with a two-stage dual-domain pipeline: standard
pretraining on a \emph{source} domain, then fine-tuning on a \emph{target} domain using the
target's own detection loss plus an MMD term aligning backbone features between the two
domains (RBF kernel, EMA-tracked bandwidth). The source domain's forward pass is
gradient-detached, so it only ever supplies a fixed feature-space reference.

\textbf{Data.} Two synthetic sources are used, both CARLA-rendered: \emph{Syndrone} (single
weather condition, vehicle annotations only) and \emph{FlyAwareV2} (four weather conditions --
clear day, clear night, rain, fog -- vehicle and person annotations). FlyAwareV2 ships only
with semantic segmentation maps, so bounding boxes were derived via a clustering plus
statistical (bell-curve) split of adjacent same-class regions -- cheaper than adapting the
original scene-generation code to emit instance boxes directly, at the cost of imperfect
instance separation for touching objects. Only the 50m-height imagery is used; 20m/80m were
omitted since the box-splitting heuristic needs height-dependent object-size priors not yet
derived.

\textbf{Regularization.} The first attempts used MMD ``as is'': a fixed-ish loss weight, and a
kernel bandwidth that is re-estimated every training step as a moving average of how far apart
the two domains' features currently are. The problem, diagnosed by visualizing the feature
space over training (via a fixed PCA basis, fit once and never refit): as the two domains'
features are pulled closer together, that bandwidth estimate shrinks too -- which sharpens the
kernel and makes it \emph{cheaper} for the optimizer to keep collapsing the features into a
tight cluster than to actually reshape the target distribution to match the source's. In other
words, the alignment mechanism was quietly rewarding a degenerate shortcut instead of genuine
alignment, and this got worse as training went on. The fix that helped: (1) \textbf{freeze the
bandwidth} after a short warm-up period, so it stops chasing the collapsing features, and (2)
\textbf{decay the MMD loss weight} linearly over training, so alignment pressure is strongest
early (before collapse can set in) and weaker later. Both together break the feedback loop; see
Section~2 for what this recovers.

\textbf{Findings so far.} In the practically relevant regime -- a large, cheap synthetic
source fine-tuned against a small real target, i.e.\ ``abundant cheap synthetic data helping
scarce real labels'' -- naive MMD showed no benefit over a plain fine-tune. Adding two
regularization mechanisms (bandwidth freeze + linear MMD-weight decay) recovered a small,
consistent target-domain benefit. Comparing the two synthetic sources under this same regime,
Syndrone (narrow, single-weather) does not clearly beat its own in-domain baseline even with
regularization, while FlyAwareV2 (diverse -- weather, occlusion, crowding) does. The working
hypothesis is that MMD's benefit depends on the synthetic source contributing genuine
additional intra-class variance beyond what the real domain already covers -- but this has not
yet been tested under matched conditions (see Section~\ref{sec:todos}).

\section{Headline result: regularization recovers a benefit (FlyAwareV2, large source / small
target)}

\begin{table}[h!]
    \centering
    \begin{tabular}{lrrrr}
        \toprule
        Model & v mAP50 & v mAP50-95 & p mAP50 & p mAP50-95 \\
        \midrule
        no-MMD (plain fine-tune control) & 0.877 & 0.626 & 0.644 & 0.303 \\
        unregularized MMD (constant weight, no freeze) & 0.877 & 0.625 & 0.642 & 0.296 \\
        regularized MMD & 0.881 & 0.628 & 0.649 & 0.300 \\
        regularized MMD, smaller decaying weight & 0.879 & 0.625 & 0.646 & 0.302 \\
        \textbf{regularized, smaller weight, 100 epochs (best)} & \textbf{0.882} & \textbf{0.630} & \textbf{0.651} & \textbf{0.304} \\
        \bottomrule
    \end{tabular}
    \caption{Validated on the real val set (target domain, the metric of primary interest).
    Unregularized MMD is statistically indistinguishable from the no-MMD control; all three
    regularized variants edge it out, and the best configuration wins on both classes.}
\end{table}

\textbf{The problem regularization fixes.} The RBF kernel's bandwidth is estimated as an
exponential moving average of the batch's mean pairwise feature distance, updated every step.
As the two domains' features are pulled together by MMD, this estimate shrinks too --
sharpening the kernel and rewarding the optimizer for collapsing the representation further
rather than genuinely matching the two distributions, a self-reinforcing feedback loop
confirmed via fixed-basis PCA visualization of the feature space. \textbf{The fix}: freeze the
bandwidth after a short warm-up period, and linearly decay the MMD loss weight over the course
of training. This breaks the loop and is what produces the table above -- without it, MMD adds
nothing in this regime.

\section{Dataset comparison: Syndrone vs.\ FlyAwareV2}

\begin{table}[h!]
    \centering
    \begin{tabular}{llrr}
        \toprule
        Model & Eval domain & mAP50 & mAP50-95 \\
        \midrule
        baseline (real, in-domain) & real & 0.896 & 0.641 \\
        syn\_source\_real\_target (unregularized) & real (target) & 0.867 & 0.599 \\
        syn\_source\_real\_target\_reg (regularized) & real (target) & 0.886 & 0.618 \\
        \midrule
        baseline (syn, in-domain) & syn & 0.758 & 0.508 \\
        syn\_source\_real\_target (unregularized) & syn (source) & 0.568 & 0.337 \\
        syn\_source\_real\_target\_reg (regularized) & syn (source) & 0.674 & 0.396 \\
        \bottomrule
    \end{tabular}
    \caption{Syndrone as source, real as target, single class (vehicle), local machine, no
    weather variation. Vehicle class only, since Syndrone has no other annotations.}
\end{table}

Regularization substantially recovers source-domain forgetting (0.568 $\to$ 0.674) and nudges
target performance up slightly (0.867 $\to$ 0.886), but \textbf{even the regularized target
result does not exceed the in-domain baseline} (0.896) -- at best, on par with training on real
data alone, not a net improvement. This is the same practically-relevant direction and regime
as Section~2's FlyAwareV2 result, which \emph{does} show a real edge over its own no-MMD
control. That contrast is the basis for the diversity hypothesis in Section~1, but see
Section~\ref{sec:todos} for why it is not yet a controlled comparison.

\section{Remaining work}
\label{sec:todos}

\subsection{A note on class mismatch}

Syndrone has vehicle annotations only; FlyAwareV2 and VisDrone both have vehicle and person.
This means two separate questions need two separate experiment designs rather than one shared
grid:

\begin{itemize}
    \item \textbf{Syndrone vs.\ FlyAwareV2 diversity comparison}: restricted to
    \textbf{vehicle-only} training and evaluation on both sides -- the only class both sources
    can supply, and it also removes a second confound (a multi-class model splits detection-head
    capacity between classes, so its vehicle mAP alone isn't directly comparable to a
    single-class model's).
    \item \textbf{Person-interference hypothesis}: inherently FlyAwareV2-only, since Syndrone
    cannot test it. Isolated by comparing a FlyAwareV2-source vehicle-only run against a
    FlyAwareV2-source multi-class (vehicle + person) run, both evaluated on real vehicle mAP,
    with batch size and regularization config matched.
\end{itemize}

\subsection{TODO 1: matched Syndrone vs.\ FlyAwareV2 comparison}

The comparison in Section~3 is confounded -- the two results were produced with different
regularization settings, possibly different batch sizes, at different points in the project.
Needed: YOLOv10n pretrained on (1) VisDrone (done), (2) Syndrone, (3) FlyAwareV2 (done), then
the same large-source/small-target fine-tune sweep for both synthetic sources at one fixed
batch size, with MMD hyperparameters logged this time rather than recalled after the fact:
no-MMD control, naive (unregularized) MMD, regularized MMD (weight config 1), regularized MMD
(weight config 2, smaller weight).

\subsection{Hypothesis 1: the person class interferes with vehicle performance}

Across the regularized variants in Section~2, person mAP consistently drops while vehicle mAP
rises, and the ``best'' model is currently selected on combined performance across both
classes. It's possible the small, blurry synthetic person instances are dragging down model
selection in a way that masks a cleaner vehicle-only signal. Test per the class-mismatch note
above: FlyAwareV2-source vehicle-only vs.\ multi-class, same regularization config, compare
vehicle mAP.

\subsection{Hypothesis 2: push the synthetic:real ratio further}

The current large-source/small-target setup is already asymmetric (roughly 10k synthetic to
1.6k real, $\sim$6:1). To more aggressively test the paper's core premise -- that an abundant,
cheap synthetic pool can compensate for scarce real labels -- push this ratio further toward
$\sim$10:1 and see whether the target-domain benefit strengthens, weakens, or holds steady.

\subsection{Final planned experiment}

Vehicle-only, at full scale ($\sim$10{,}000 synthetic FlyAwareV2 images, $\sim$1{,}500 real
VisDrone images), on both YOLOv10n and YOLOv10m, each with the four fine-tune variants (no-MMD,
naive MMD, regularized MMD $\times$2 weight configs). This is the result intended to close out
the paper's central claim.

\section{Related work}

\begin{itemize}
    \item \textbf{Long et al., Deep Adaptation Network (DAN) / Multi-Kernel MMD} -- the
    original MMD-in-a-deep-net domain adaptation method, and specifically the multi-kernel
    variant designed to reduce sensitivity to kernel bandwidth choice. Directly relevant: this
    is the standard answer to the same brittleness our bandwidth-freeze regularization
    addresses differently.
    \item \textbf{Tzeng et al., Deep Domain Confusion} -- earlier precursor to DAN, establishes
    the MMD-as-a-training-loss lineage this paper sits in.
    \item \textbf{Long et al., Joint Adaptation Networks (JAN)} -- extends MMD alignment to
    joint distributions across layers; worth citing as a more sophisticated alternative not
    pursued here.
    \item \textbf{Sun \& Saenko, CORAL} -- second-order-statistics alternative to MMD;
    contrasts our discrepancy-minimization choice against a simpler covariance-matching one.
    \item \textbf{Ganin \& Lempitsky, DANN} -- the adversarial (gradient-reversal) alternative
    to discrepancy-based alignment; needed to place MMD among the two dominant families of
    feature-alignment domain adaptation.
    \item \textbf{Chen et al., Domain Adaptive Faster R-CNN} -- adversarial domain adaptation
    applied specifically to object detection, the closest prior work in task (detection) even
    though it differs in alignment mechanism (adversarial vs.\ MMD).
    \item \textbf{Richter et al., Playing for Data (GTA5); SYNTHIA} -- canonical synthetic
    driving-scene datasets for perception, establishing the sim-to-real-for-perception lineage
    this paper extends to the aerial setting.
    \item \textbf{Dosovitskiy et al., CARLA} -- the simulator underlying both Syndrone and
    FlyAwareV2; cited directly as the data source.
    \item \textbf{Tobin et al., domain randomization} -- an alternative synthetic-data strategy
    (randomize simulation parameters rather than align features post-hoc) not used here; useful
    as a contrasting future-work direction.
    \item \textbf{Zhu et al., VisDrone} -- the real-domain benchmark and data source.
    \item \textbf{UAVDT, DOTA} -- sibling aerial-imagery detection benchmarks, cited briefly to
    place VisDrone among them.
    \item \textbf{Akyon et al., SAHI} -- the standard tiling-based alternative to large input
    resolution for small-object detection in aerial imagery; relevant to justify why this work
    uses a large \texttt{imgsz} instead of slicing.
    \item \textbf{Syndrone and FlyAwareV2 original papers} -- cited directly as the data
    sources; worth checking whether their authors already report a sim-to-real gap on the same
    data, as a number to compare against.
    \item \textbf{Connected-components / watershed / clustering-based instance separation
    literature} -- situates the segmentation-to-bounding-box conversion method (used for
    FlyAwareV2) as an instance of a known problem rather than an ad hoc heuristic.
\end{itemize}

\end{document}
